% Introduction (about 1.5 pages in a two-column conference format).
\section{Introduction}
\label{sec:intro}

Visual document understanding (VDU) is the task of reading a document the way a person does and getting useful information from it. A document is more than text. It also has tables, charts, figures, images and a page layout. A VDU system must read all of these together \cite{xu2020layoutlm,kim2022donut}. Early methods ran optical character recognition (OCR) first and then used the text and its position on the page \cite{xu2020layoutlm}. Later methods read the page image directly \cite{kim2022donut}. Today, large vision-language models (VLMs) such as Qwen3-VL \cite{bai2025qwen3vl} take page images as input and can answer questions about them.

A main VDU task is visual question answering (VQA). In VQA, the system gets an image and a question in natural language, and it must return a correct answer \cite{antol2015vqa}. When the image is a document page, the task is called document VQA (DocVQA) \cite{mathew2021docvqa}. For example, a question may ask for the total in a table of an invoice, or for the highest bar in a chart. The answer is usually a short text, a number or a list. The first DocVQA datasets use one page image per question \cite{mathew2021docvqa}. This is much easier than the real use case, where the evidence is spread over many pages.

Document question answering matters because people need facts from documents every day. A worker may need one number from a financial report. A student may need one result from a research paper. A technician may need one step from a long manual. A clerk may need a date or an amount from an official form, and a lawyer may need one clause from a long contract. Reading all the pages takes a lot of time, so a good system saves time and money. A wrong answer can also cause harm, for example a wrong amount in a financial report or a wrong step in a repair guide. So the system must be both fast and correct.

Real documents are often long. Reports, manuals and research papers have dozens of pages, and the evidence for one question can be in a table on one page and in a chart on another. MMLongBench-Doc \cite{ma2024mmlongbench} is a benchmark for this setting. Its documents have 49.4 pages on average, about one third of its questions need evidence from more than one page, and about 23\% of its questions cannot be answered from the document. In the benchmark paper, the best model, GPT-4o, reaches only 42.7 F1, and the best open-source VLM is 11.9\% behind the best language model that reads the OCR text \cite{ma2024mmlongbench}. A long document is hard for two reasons. The system must first find the pages and the parts that hold the evidence, and then it must use that evidence to give the right answer.

A single model that reads all the pages at once has two problems. The input is very long, and small details in charts and tables are easily lost. For this reason, many recent systems use several steps or several agents for this task. DocLens \cite{zhu2025doclens}, SlideAgent \cite{slideagent}, Doc-V* \cite{docvstar} and ORCA \cite{orca} are four examples. DocLens has two modules. The Lens Module finds the right pages and the right parts of the pages. In the Reasoning Module, an Answer Sampler writes several candidate answers, and an Adjudicator then decides the final answer. DocLens reports an accuracy of 61.1 with Qwen3-VL-8B on 200 random questions of MMLongBench-Doc \cite{zhu2025doclens}. SlideAgent turns on its specialized agents only when they are needed \cite{slideagent}, and ORCA calls its debate and adjudication step only when it is needed \cite{orca}. % VERIFY: the original draft said that ORCA and SlideAgent skip later steps when agents agree, and that ORCA reports 12% of failures from over-refinement. I could not confirm this from the abstracts, so the sentence above uses only what the abstracts state.

In DocLens, the Adjudicator runs on every question. This adds one more model call to each question. The Adjudicator can also change an answer that was already correct. We call this over-correction. In this case, the Adjudicator adds cost and lowers accuracy at the same time. A second problem is that the Adjudicator reads only the text of the candidate analyses and answers. It does not see the document pages. When the candidates disagree, it cannot check them against the pages. So the main questions are when the Adjudicator is needed, and how it can use the pages when it is needed. This is the gap that we study. If the answers of the Answer Sampler agree, an extra model call brings little new information and may only add cost or risk. If they disagree, the text of the answers alone is often not enough to find the correct one, and the pages are the best source of evidence. This gives two ideas. First, we should call the Adjudicator only when the candidate answers disagree. Second, when we do call it, it should look at the pages.

To address this problem, we propose a confidence gate for DocLens. The gate counts how many of the $k=8$ candidate answers agree. If all of them agree, the system skips the Adjudicator and returns the majority answer. If they do not agree, a visual Adjudicator runs. It reads the same evidence pages as the Answer Sampler, together with the candidate answers, and then picks the final answer. The gate does not change the DocLens prompts for the Answer Sampler, and we do not fine-tune the backbone model.

\begin{enumerate}
  \item We propose an agreement gate for the DocLens Adjudicator. It normalizes the $k=8$ sampled answers and skips adjudication when all of them agree, which removes the Adjudicator call for these questions.
  \item We propose a visual Adjudicator. It conditions on the evidence pages (page screenshots, localized elements and OCR text) in addition to the candidate analyses and answers, and it is called only when the sampled answers disagree.
  \item We compare the fixed unanimity gate with a threshold chosen by five-fold cross-validation over documents, with always running the Adjudicator, and with majority voting. We report accuracy, over-correction rate, fix rate, Adjudicator rate, token cost and paired bootstrap confidence intervals.
  \item We reproduce vanilla Qwen3-VL-8B, Qwen3-VL-8B with OCR and DocLens on a random sample of 200 MMLongBench-Doc questions (seed 0)\ifextended, and we evaluate our method on 400 further questions (seed 1) that were not used to design it\fi. We explain how our numbers differ from the original paper.
\end{enumerate}
